\documentclass[../main.tex]{subfiles}
\begin{document}
\section{テスト実験データを用いたハイレベルフィルタの性能評価}
\label{sec:hlf-evaluation}
\begin{tcolorbox}
[
  title=執筆メモ,
  colback=yellow!8,
  colframe=black!30,
  breakable
]
\begin{itemize}
  \item 目的
  \item 評価方法
  \begin{itemize}
    \item データリプレイヤー
    \item T103のデータ
    \begin{itemize}
      \item T103とは
      \item ビーム
      \item セットアップ
      \begin{itemize}
        \item 検出器の配置
        \item 設置座標(おもにzが重要)
        \item UTOF, LTOF, BFT, KLDCは必須で他は確認できれば入れておく
      \end{itemize}
    \end{itemize}
    \item 評価指標
  \end{itemize}
  \item 実装して評価したハイレベルフィルタ3種
  \begin{itemize}
    \item TOFフィルタ
    \item クラスタリングフィルタ
    \item トラッキングフィルタ
  \end{itemize}
\end{itemize}
\end{tcolorbox}
MARQスペクトロメータで用いることを想定したハイレベルフィルタを複数実装し、連続読み出し式DAQテスト実験T103で取得したデータを用いてハイレベルフィルタの性能評価を行った。


\subsection{評価方法}
\begin{tcolorbox}
[
  title=執筆メモ,
  colback=yellow!8,
  colframe=black!30,
  breakable
]
\begin{itemize}
  \item 修論として結果をまとめるなら、計測環境がバラバラなのはだめだ(run000400というパラメータ調整中のデータでの計測であったことも改善すべき点の一つ。だし、クラスタリングの時は違うランを使っていたような気がしなくなくもない)
  \item 世代が一番新しくて、メモリ転送速度も現代的なmarq07で統一して計測するのが今後の見通しが良さそう
  \item それぞれの環境でなんらかのベンチマークをして計算能力の相互変換をできるようにするものいいと思う
  \begin{itemize}
    \item marq07で全部のベンチマークをする
    \item データは全部marq08(ストレージたくさん)に置いて、marq08からリプレイヤーで飛ばす
    \item e50server03は大阪に送ってもらって、すでにセットアップ済
  \end{itemize}
  \item ここ(\lstinline|/Users/kashima/codex-cli-workspace/master-thesis/workspace/20261112/common-benchmark|)でどこに何のコードがあるかまとめることにする
  \begin{itemize}
    \item TOF: 計測システム研究会2025 (ProcessSlice定義: \path{/Users/kashima/codex-cli-workspace/master-thesis/workspace/20261112/common-benchmark/src/TOF-20251111/nestdaq-user-impl_MARQ/FilterTimeFrameSliceBySomething.cxx}:177)
    \begin{itemize}
      \item リポ: \url{https://github.com/SATSUKIUM/nestdaq-user-impl_MARQ}
      \item コミットハッシュ: \texttt{53d8724486561cf40ce6de6d2c1bbd67093a9054}（chmap導入前の基準版）
      \item チェック項目
      \item ConditionalRun()のスループット(送受信は除く)
      \item その分布
      \item ProcessSlice()だけのスループット
      \item その分布
      \item $\uparrow$ の横軸データサイズ縦軸処理時間の分布
      \item カット前TOF分布図(ptag132.pdfと矛盾はないか)
      \item カット範囲の決定(ptag132.pdfと矛盾はないか)
      \item カット前後のデータ削減率(ptag132.pdfと矛盾はないか)
      \item カット後TOF分布図
    \end{itemize}
    \item Multiplicity: 2026春物理学会 (ProcessSlice定義: \path{/Users/kashima/codex-cli-workspace/master-thesis/workspace/20261112/common-benchmark/src/Multiplicity/nestdaq-user-impl_MARQ/FilterTimeFrameSliceByMultiplicity.cxx}:112)
    \begin{itemize}
      \item リポ: \url{https://github.com/SATSUKIUM/nestdaq-user-impl_MARQ}
      \item コミットハッシュ: \texttt{b9f8553d5524fa5f845cfa8160f06ee2da5026a8}（基準コミット。ローカル版は最新channel-map対応・配列境界修正あり。2026春発表版との一致は未確認）
      \item チェック項目
      \item ConditionalRun()のスループット(送受信は除く)
      \item その分布
      \item ProcessSlice()だけのスループット
      \item その分布
      \item $\uparrow$ の横軸データサイズ縦軸処理時間の分布
      \item カット前クラスタ数分布図
      \item カット条件ごとのデータ削減率
      \item カット前後でのオフライン解析のトラック数の変化
    \end{itemize}
    \item Track: 2026夏物理学会 (ProcessSlice定義: \path{/Users/kashima/scp_landing_zone/20260519/nestdaq-user-impl/FilterTimeFrameSliceByTrack.cxx}:452)
    \begin{itemize}
      \item リポ: \url{https://github.com/SATSUKIUM/nestdaq-user-impl}
      \item コミットハッシュ: \texttt{0d7bbf643c53e990a65c92ae41c5cb3b70de33a9}（参照版。2026夏発表版との一致・共通ベンチマークへの採用は未確定）
      \item チェック項目
      \item ConditionalRun()のスループット(送受信は除く)
      \item その分布
      \item ProcessSlice()だけのスループット
      \item その分布
      \item $\uparrow$ の横軸データサイズ縦軸処理時間の分布
      \item ビームプロファイル(BFTを使って)
      \item トラッキング結果(トラックのヒットレイヤ数, 残差分布, ヒットプロファイル, $\chi^2$分布)
      \item レイヤごとの検出効率(BFTのトラッキング結果と比較)
    \end{itemize}
  \end{itemize}
\end{itemize}
\end{tcolorbox}
記録済みのT103データをリプレイヤーから入力し、1プロセスのオンラインフィルタが処理したデータ量と処理時間を測定した。
TOFフィルタとクラスター数フィルタの評価にはAMD EPYC 74F3、64\,GBメモリの計算機を用いた。
トラッキングフィルタの評価にはAMD Ryzen 9 7900、128\,GBメモリ、AlmaLinux 9.8の計算機を用い、省電力機能を無効化した\cite{kashima:tof2025,kashima:cluster2026,kashima:tracking2026}。
以下の処理性能は、spill on中のフィルタ全体について求めたスループット分布の平均値である。
データの展開や送受信を含む全体処理と、判定関数のみの処理時間を区別する。
ここではMBを$10^6$\,byteとする。

\subsection{TOFフィルタ}

タイミングカウンタのヒット時刻の差から飛行時間を求め、その範囲を条件として事象を選別する。
T103データでは飛行時間の分布を用いて$K^{-}$ビームの事象を選別した。
% TODO: 元TOFヒストグラムを確認後、05-hlf/tof-selectionで再作図する。


フィルタ全体の平均スループットは594.9\,MB/sであった\cite{kashima:tof2025}。
CPUの省電力・ブースト機能を無効化し、全コアを3.2\,GHzに固定した評価である。
TOF計算・判定だけでなく、受信データの展開や送信などを含むフィルタ全体の値を用いた。
ヒット数が増えるとカウンタ間の組み合わせも増えるため、高ビーム強度への適用では処理性能の入力密度依存性を考慮する必要がある。

\begin{tcolorbox}[
  title=執筆メモ,
  colback=yellow!8,
  colframe=black!30,
  breakable
]
\begin{itemize}
  \item 他のハイレベルフィルタの基礎となる共通の処理
  \item TOFを計算することで、粒子識別のために重要である粒子の速度を求めることができる
  \begin{itemize}
    \item 後で確認
    \item 例えばLTOFだとセグメント数は効かない
    \item むしろイベントウィンドウの中のビーム密度の違い
  \end{itemize}
\end{itemize}
\end{tcolorbox}

\subsection{クラスタリングフィルタ}

KLDCのヒットをレイヤごとにクラスタリングし、クラスター数が2以上のレイヤが$N$枚以上であることを選別条件とした。
UとU$'$、VとV$'$のペアをそれぞれ同一視し、2台のKLDCを計4レイヤとして扱った。
対象はコインシデンス時刻から150\,ns以内のヒットであり、1ワイヤのヒットからクラスターとして数えた\cite{kashima:cluster2026}。

$N=2$の条件では、全エントリ数は387974から168213へ減少した。
オフライン解析で再構成トラック数が2以上となる事象の割合は約5.6\%から約10.8\%へ増加した。
一方で、その件数は21659から18118へ減少している。
この割合の増加は選別後の試料中の構成比の変化であり、目的事象の保持効率を表すものではない。
図\ref{fig:hlf-cluster-selection}に選別条件とエントリ数の関係を示す。
フィルタ全体の平均スループットは442.9\,MB/sであった\cite{kashima:cluster2026}。

\begin{figure}[htbp]
  \centering
  \makebox[\linewidth][c]{\includegraphics[pagebox=cropbox,width=0.95\linewidth,height=0.70\textheight,keepaspectratio]{05-hlf/cluster-selection/cluster-selection.pdf}}
  \caption{クラスター数による選別とオフライン再構成トラック数の関係。$N$はクラスター数が2以上のレイヤに課す最小枚数\cite{kashima:cluster2026}}
  \label{fig:hlf-cluster-selection}
\end{figure}

\begin{tcolorbox}[
  title=執筆メモ,
  colback=yellow!8,
  colframe=black!30,
  breakable
]
\begin{itemize}
  \item クラスタリングを行うことで、イベントが解析可能であるための最小の条件を課す。少ない計算量でもフィルタリング後のイベントの解析可能性の向上が期待できる
\end{itemize}
\end{tcolorbox}

\subsection{トラッキングフィルタ}

ドリフトチェンバーのTDCデータをヒット情報に変換し、ペア面の隣接ワイヤヒットをクラスタリングする。
ヒット位置と直線飛跡との残差に対する最小二乗フィッティングにより、飛跡の切片$(x_0,y_0)$と傾き$(dx/dz,dy/dz)$を求める。
再構成した位置と傾きは、ビーム粒子と散乱粒子を判別するための情報となる\cite{kashima:tracking2026}。

フィルタ全体の平均スループットは123.4\,MB/sであった。
処理時間の大部分は飛跡候補の生成とフィッティングに費やされた\cite{kashima:tracking2026}。
クラスター数による選別と比べて詳細な飛跡情報を得られるが、計算負荷も大きい。
両フィルタの性能測定ではCPUが異なるため、処理時間の差をアルゴリズムの違いだけに帰することはできない。


\begin{tcolorbox}
[
  title=執筆メモ,
  colback=yellow!8,
  colframe=black!30,
  breakable
]
\begin{itemize}
  \item ローカル直線トラッキングを行うことで、トラッカーに残された粒子の飛跡へアクセスできる。
  \begin{itemize}
    \item 後で確認
    \item DCの面が変わった時の処理速度の違い
  \end{itemize}
\end{itemize}
\end{tcolorbox}


\subsection{処理性能のまとめ}
表\ref{tab:hlf-performance}に各フィルタの処理性能を示す。
いずれも1プロセスあたりの評価である。
\begin{table}[htbp]
  \centering
  \caption{spill on中のオンラインフィルタ全体の平均スループット}
  \label{tab:hlf-performance}
  \begin{tabular}{lll}
    \hline
    フィルタ & CPU & スループット [MB/s] \\
    \hline
    TOF & EPYC 74F3 & 594.9 \\
    クラスター数 & EPYC 74F3 & 442.9 \\
    ローカルトラッキング & Ryzen 9 7900 & 123.4 \\
    \hline
  \end{tabular}
\end{table}

\end{document}
